\section{Conclusion}\label{sec:conclusion}

We tested Theorem~3 of Operon's external-frameworks integration
memo --- the conjecture that certificate-based evaluators converge
faster than scalar-reward evaluators on tasks with enumerable
failure modes --- by running GEPA with three evaluator shapes on a
synthetic windowed-stability task.  The cert-binary arm replaces
the scalar reward with a pass/fail verdict plus structured
obligation text drawn from an Operon theorem verifier.  The
scalar-reward and scalar-with-evidence arms form a minimal pair
that isolates binarization from obligation-text content.

\textbf{Verdict: refine.}  All three arms reach 100\% convergence
within budget; mean iterations-to-convergence are 1.0
(cert-binary), 2.0 (scalar), and 1.0 (scalar-evidence).  Against
the stronger baseline (scalar-evidence) cert-binary is statistically
indistinguishable (Mann-Whitney $U=50$, $p=1.000$, rank-biserial
$0.00$), so the predeclared $\le 0.75$-ratio gate is not cleared
and Theorem~3 stays a conjecture rather than an accepted theorem
for this task.  The formatter ablation on seeds 0--4 ties at 1.0
mean iterations for both the default and the minimal formatter,
localising the active ingredient to ``\emph{some} structured
failure information vs none'' --- neither binarization, theorem
framing, nor the numeric content of the obligation evidence adds
on top.  The integration memo is updated accordingly: certificates
remain a clean evaluator interface, but on this task they bring no
convergence-speed advantage over a graded scalar reward paired with
the same evidence text.

\subsection{Follow-ups unblocked}\label{sec:followups}

\begin{enumerate}[leftmargin=*,itemsep=2pt,topsep=2pt]
\item \textbf{Theorem-shape sweep}: rerun the protocol on
\texttt{behavioral\_quality} (continuous) and
\texttt{behavioral\_no\_anomaly} (structural).
\item \textbf{Cross-LM robustness}: rerun cert-binary and the
strongest baseline with a frontier LM.
\item \textbf{Operon--autocontext adapter}: wire Operon certificates
into autocontext's staged-validation loop~\cite{autocontext2025},
giving a second external framework that consumes certificates as
evaluator output.  A plan for this is logged separately; this is a
paper-7 candidate.
\item \textbf{Lightweight DSPy reproducibility theorem}: ship
\texttt{dspy\_compile\_pinned\_inputs} as the cheap variant of
memo Theorem~2.
\item \textbf{\texttt{gepa\_candidate\_improvement} theorem}:
certify that an evolved candidate Pareto-dominates its parent on
the validation set; close GEPA's outer loop with a certificate.
\end{enumerate}

\subsection{Artifacts}

\begin{itemize}[leftmargin=*,itemsep=2pt,topsep=2pt]
\item Experiment driver:
\texttt{eval/convergence/theorem\_6\_experiment.py}
\item Harness:
\texttt{eval/convergence/synthetic\_signal\_harness.py}
\item Baseline adapters:
\texttt{scalar\_reward\_adapter.py},
\texttt{scalar\_with\_evidence\_adapter.py}
\item Analysis:
\texttt{eval/convergence/analysis\_theorem\_6.py}
\item Per-arm trajectories:
\texttt{eval/results/theorem\_6/\{arm\}/seed\_*.json}
\item Summary:
\texttt{eval/results/theorem\_6/summary.json}
\item Unit tests:
\texttt{tests/unit/convergence/test\_paper6\_*.py}
\end{itemize}

\paragraph{Reproducibility.}
\begin{verbatim}
pip install gepa matplotlib scipy
ollama pull gemma4  # if not already local
python -m eval.convergence.theorem_6_experiment --sweep
python -m eval.convergence.analysis_theorem_6
cd article/paper6 && tectonic main.tex
\end{verbatim}
